\documentclass{situationist-preprint}

\title{The Shipowner Case}
\msid{EPI-23-b}
\msfamily{EPI-23}
\msdiscipline{Epistemology}
\mstype{Companion manuscript}
\mskeywords{ethics of belief; Clifford; James; shipowner; culpable inquiry; moral luck}

\begin{document}
\maketitle

\begin{abstract}
Clifford's shipowner suppresses his doubts about an old ship, sincerely comes to believe she is seaworthy, and sends her to sea with emigrants aboard. Clifford held that he is guilty, and would be guilty even had the ship arrived safely. We show that the account of EPI-23-a explains both verdicts: the shipowner's case falls squarely within Clifford's domain, where the stakes of error are high and evidence is available, and the wrong lies in the failure of inquiry, which is the same whether the ship sinks or arrives. Varying the case shows which features carry the verdict.
\end{abstract}

\section{The case}

The shipowner knows that his ship is old and has often needed repairs; he is advised to have her overhauled. He reflects that she has come through many voyages, puts his doubts aside, and comes to believe sincerely that she is safe. She sinks in mid-ocean \citep{Clifford1877}.

By Definition~2.2 of EPI-23-a the stakes of error are high: if the ship is unseaworthy, sailing risks the lives of the passengers, and withholding judgment and delaying the voyage would avoid that risk. By Definition~2.3 evidence is available: an overhaul would settle the question, at a cost proportionate to the stakes, before the ship must sail. So Theorem~3.1 of EPI-23-a applies: believing on insufficient evidence and acting on the belief is a wrong.

\begin{proposition}[Outcome-independence]\label{p:outcome}
The shipowner does wrong whether or not the ship arrives safely.
\end{proposition}

\begin{proof}[Argument]
The conditions of Theorem~3.1 of EPI-23-a concern the evidence, the stakes and the availability of inquiry at the time of belief and action, and the wrong it identifies is the culpable failure to inquire. None of these depends on the outcome of the voyage.
\end{proof}

This is Clifford's own second verdict: a safe arrival would not have made the shipowner innocent, only undiscovered. The account explains why, without denying that the sinking adds a further harm for which he is also answerable.

\section{Variations}

\emph{No passengers.} Suppose the ship is to carry only cargo the owner himself has insured and is willing to lose, and an overhaul is available. The stakes of error are not high in the sense of Definition~2.2, and Theorem~3.1 does not apply. His belief is still epistemically unfit, by Principle~P1 of EPI-23-a, but Clifford's charge of wrongdoing loses its ground.

\emph{No inspection possible.} Suppose a storm arises at sea and the captain must decide whether the hull will hold if he runs before it. The stakes are high, but no further evidence can be had before he must act. Neither domain of EPI-23-a contains this case; it is decided by acting on the rational credence, by Principle~P2 there, and we do not extend either principle to it.

\begin{objection}
The shipowner is sincere; one cannot be blamed for a belief one honestly holds \citep[cf.][]{James1896}.
\end{objection}

\begin{reply}
Sincerity concerns whether he believes; the wrong concerns how he came to. He suppressed doubts that called for inquiry he could have made. That is a failure for which he is answerable, even though belief is not under direct voluntary control, because the inquiry was.
\end{reply}

\section{Conclusion}

The shipowner case is the paradigm of Clifford's domain. It supports Clifford's principle exactly where the stakes are high and evidence is available, and the variations show that removing either condition removes the case's force.

\begin{thebibliography}{99}\small
\bibitem[Clifford(1877)]{Clifford1877} Clifford, W.~K. (1877). The ethics of belief. \emph{Contemporary Review}, 29, 289--309.
\bibitem[James(1896)]{James1896} James, W. (1896). The will to believe. \emph{The New World}, 5, 327--347.
\end{thebibliography}

\end{document}
